\documentclass[../main.tex]{subfiles}
\begin{document}
\section{Tier 4 --- heavy tails and the LSI/\texorpdfstring{$T_2$}{T2} obstruction}
\label{sec:tier4}

Heavy-tailed priors --- Laplace, Student-$t$, Cauchy, horseshoe, scale mixtures of
Gaussians --- are statistically attractive (robust, sparsity-promoting) but analytically
hard. The single mechanism behind every difficulty in this tier is that a log-Sobolev
inequality forces sub-Gaussian concentration.

\begin{theorem}[Heavy tails forbid LSI]
\label{thm:heavy-tail-no-lsi}
If $\Pstar$ satisfies LSI with constant $\CLS$, then every $1$-Lipschitz $f$ is sub-Gaussian:
$\Pstar(f-\E f\ge r)\le\exp(-r^2/2\CLS)$. Consequently a target with merely exponential or
polynomial tails satisfies \emph{no} classical LSI, and (being tied to Gaussian
concentration) no $T_2$ either.
\end{theorem}
\noindent\cite{BobkovGotze1999,OttoVillani2000}. The escape is to state results at the
weak-/weighted-Poincar\'e level, or to reparameterize to sub-Gaussian coordinates before
certifying.

\begin{family}[Classification of heavy-tailed priors by the strongest inequality they satisfy]
\label{fam:heavy-tail-classification}
Under the Euclidean gradient form,
\begin{center}
\begin{tabularx}{\textwidth}{@{}p{0.26\textwidth} c c c Y@{}}
\toprule
Prior / target & Poincar\'e & LSI & $T_2$ & Reason\\
\midrule
Gaussian                 & yes & yes & yes & sub-Gaussian tails\\
Laplace $\propto e^{-\abs{x}/b}$ & $\CP\sim b^2$ & \textbf{no} & \textbf{no} & exponential tails\\
Gamma                    & usually & \textbf{no} & \textbf{no} & exponential tail\\
Student-$t_\nu$          & \textbf{no} / weighted & \textbf{no} & \textbf{no} & polynomial tail\\
Cauchy                   & no / weak & \textbf{no} & \textbf{no} & too heavy-tailed\\
horseshoe                & delicate & \textbf{no} & \textbf{no} & scale-mixture tail\\
\bottomrule
\end{tabularx}
\end{center}
\end{family}

\noindent The Laplace law $\Pstar(\dd x)=\tfrac1{2b}e^{-\abs{x}/b}\dd x$ has a finite
Poincar\'e constant $\CP\sim b^2$ (Hardy criterion, Theorem~\ref{thm:hardy-1d}) but only
exponential concentration, so no LSI. For Student-$t_\nu$ the tails are polynomial: LSI and
$T_2$ fail, and the \emph{classical} Poincar\'e inequality fails as well for \emph{every}
$\nu$, not only small $\nu$, because a polynomial tail makes the $1$D Hardy functional diverge.
The correct object is a \emph{weighted} Poincar\'e inequality, finite for all $\nu$ but with a
constant that worsens as the tail heavies; Cauchy ($\nu=1$) is the heaviest case
(Section~\ref{sec:a3} develops the quantitative weighted theory). The right replacement tools
are the modified / weighted log-Sobolev inequalities and the $1$D Hardy theory; making this
classification quantitative for hierarchical robust-Bayes posteriors is part of the agenda
(Section~\ref{sec:agenda}).
\end{document}
